% !TeX encoding = UTF-8
% !TeX program = xelatex


\documentclass[12pt,openright]{book}





\usepackage[]{NKThesis}   % 中文
%\usepackage[English]{NKThesis} % 英文



\usepackage{amssymb}


\usepackage{amsmath}
\usepackage{cite}

\graphicspath{{fig/}}

% 为了分段编译更快
\includeonly{
./tex/abstract,
./tex/manual,
./tex/acknowl,
./tex/resume
}

\newtheorem{Theorem}{\hskip 2em 定理}[chapter]
\newtheorem{Lemma}[Theorem]{\hskip 2em 引理}
\newtheorem{Corollary}[Theorem]{\hskip 2em 推论}
\newtheorem{Proposition}[Theorem]{\hskip 2em 命题}
\newtheorem{Definition}[Theorem]{\hskip 2em 定义}
\newtheorem{Example}[Theorem]{\hskip 2em 例}



\begin{document}

%  设置基本信息
%  注意:  逗号`,'是项目分隔符. 如果某一项的值出现逗号, 应放在花括号内, 如 {,}
%
\NKTsetup{%
  论文题目(中文) = LaTeX论文模板,
  副标题         = ——论文副标题,
  论文题目(英文) =  A LaTeX thesis template,
  论文作者       = 张十三,
  学号           = 1020161111,
  指导教师       = 张三\quad 教授,
  申请学位       = 理学硕士,
  培养单位       = 某某学院,
  学科专业       = 某某某某,
  研究方向       = 某某某某,
  中图分类号     = ,
  UDC            = ,
  学校代码       = 10055,
  密级           = 公开,
                   % 公开 | 限制 | 秘密 | 机密, 若为公开, 不填以下三项
  保密期限       = ,
  审批表编号     = ,
  批准日期       = ,
  论文完成时间   = 二〇一九年三月,
  答辩日期       = ,
  论文类别       = 硕士专业学位,
                   % 博士 | 专业博士 | 学历硕士 | 硕士专业学位 | 本科
  答辩委员会主席       = 李四,
  评阅人       = 李四\quad 王五,      
  专业           = 某某某某,
  联系电话       = 1234567,
  Email          = 123@123.com,
  通讯地址(邮编) = 天津市南开区卫津路94号(300071),
  备注           = {},
}



% -*- coding: utf-8 -*-


\begin{zhaiyao}
自从1998年两个独立的超新星研究小组分别发现宇宙正在加速膨胀，暗能量就成为宇宙学中的热门。
\end{zhaiyao}




\begin{guanjianci}
暗能量；动力学分析
\end{guanjianci}



\begin{abstract}
Since 1998, two independent supernova research groups have discovered that the universe is accelerating, and the dark energy has become a hot topic in cosmology. 

\end{abstract}



\begin{keywords}
dark energy model; dynamic analysis
\end{keywords} 


% -*- coding: utf-8 -*-

\chapter{南开大学学位论文格式宏包 NKThesis 使用说明} \label{chpt:A}

\section{系统要求}

模板在 TexLive 2020/2021/2022 下测试\XeLaTeX, PDF\LaTeX 通过。
对于其它 TeX 发行版可能需要做个别修改。


本模板最初版来自孙老师\url{http://222.30.48.141/$\sim$sunwch}的版本，
格式参照南开大学2019年毕业论文格式。
后来在志愿者的努力下做成了开源项目
\url{https://github.com/NewFuture/NKThesis}。
%
同时，赵子懿，潘宇等同学也主动维护了本科论文的\LaTeX 开源版本
\url{https://www.overleaf.com/latex/templates/nku-dissertation-template2020/ryscpwbhbmhs}。
%
鉴于各种版本维护复杂，2022年12月疫情期间，南开大学计算机系程明明教授
在硕士、博士论文的基础上，统一了本科、学术硕士、专业硕士、博士和专业博士的
论文模板。
新的模板在GitHb上开放共享\url{https://github.com/MCG-NKU/NKU-Thesis}，
欢迎有问题的同学及时反馈。

\section{NKThesis 使用说明}

本模板可以使用以下两种方式编译:
\begin{enumerate}
 \item \XeLaTeX + BibTex [推荐]
 \item PDF\LaTeX + BibTex
\end{enumerate}


我们建议您使用\XeLaTeX 编译。与PDF\LaTeX 相比，\XeLaTeX  编译长文档的速度更快，
编译一篇一百多页的论文只需几秒的时间(SL9400 @ 1.86GHz)。

在改变编译方式前应先删除 *.toc 和 *.aux 文件，
因为不同编译方式产生的辅助文件格式可能并不相同。


在Windows系统中可以直接使用PDF\LaTeX。
\textcolor{red}{如果想在Overleaf中使用Pdf\LaTeX，
需要传相关字体文件到根目录：
simfang.ttf，SimSun.ttc，simkai.ttf，SimHei.ttf。}


本模板用到 宋体、楷体、仿宋、黑体四种字体. 若需重新配置字体, 请自行添加修改。%请修改 NKTfonts.cfg。
对于 Linux/Mac, 可能需要设置环境变量 OSFONTDIR, 具体内容请参考 texmf.cnf。



\section{引用章节号}
\label{sec:ex:A}

引用章节号请参考如下格式: \ref{chpt:A}\ref{sec:ex:A}.


\section{中英文间隔}

使用 \XeLaTeX\ 编译时，会自动在中英文转换时添加必要的空格。 使用 [PDF]\LaTeX\
编译时仅忽略中文之间的空格，而中英文之间的空格予以保留。
因此，不管何种编译方式，您都不需要在中英文间添加 $\tilde{}$ 以获得额外的空格。例如，

这是 English 中文 $x=y$ 测试

这是English中文$x=y$测试

可以看出，以上两行用 \XeLaTeX\ 编译的结果是相同的。


\section{NKThesis 预调用的宏包}

NKThesis 已经调用以下宏包，您无须重新调用。

\begin{center}
\tablecaption{NKThesis 预调用的宏包}
\begin{tabular}{l|l}
\hline
编译方式 & 调用的宏包\\ \hline
\XeLaTeX & xeCJK, CJKnumb, graphicx, mathptmx \\ \hline
[PDF]\LaTeX & CJK, CJKnumb, uniGBK, graphicx, mathptmx \\
\hline
\end{tabular}
\end{center}


\section{图表}

插图的例子：

\begin{center}
\includegraphics[viewport=0 0 2984 969,width=40mm]{nankaidaxue.pdf}
\figurecaption{南开大学}
\end{center}

\section{字体}

一般情况下, 您不需要显式地设置字体. 如果确实需要, 请使用以下命令

\begin{verbatim}
宋体:  \rmfamily\upshape 或 \songti
黑体:  \bfseries 或 \heiti
楷体:  \itshape  或 \kaiti
仿宋:  \ttfamily 或 \fangsong
加粗:  \jiacu
\end{verbatim}


\section{参考文献} \label{manual:ref}
参考文献引用:
\cite{ChenCheChen2001,Nadkarni-1992,Hua-Wang-1973}
文献\cite{ZhuKeZhen,Huo}\cite{Timoshenko,Zhang-Wang}
\cite[Theorem 2.1]{ZhuKeZhen,21PAMI-Sal100K}

\subsection{录入参考文献}

本模板采用 bibtex 宏包管理参考文献。如果你对此不熟悉，可以
\begin{enumerate}
\item 参考宏包使用说明，或者
\item 手工排版参考文献，然后参考 nkthesis.bib 最后 3 条的格式录入。
\end{enumerate}



\section{一些建议}
\subsection{关于分数的写法}


\LaTeX 提供宏命令\verb+\frac+, 用以打印分数. 为使得版面整齐, 该命令的使用应遵循以下原则:

\begin{enumerate}
\item 仅在分行表达式中使用,
\item 不嵌套使用,
\item 不在上下标中使用.
\end{enumerate}

也就是说, 行内表达式和上下标中出现分数时一律用 $a/b$表示, 如
$(x+2)/((3x^2+4)(7+y))$. 下面是居中表达式:

\[
 x^2 = y^{1/2} +3.
\]

多行表达式: 尽量在加、减、乘、等号前换行. 在乘号前换行时,
下一行首用 \string\times:
\def\iint{\mathop{\int\!\!\!\int}}\def\calG{\mathcal G}
\begin{eqnarray}
&&\left|(W_{\psi_1}f)(a,b)-(W_{\psi_1}f)(a_j,b_{j,k})\right|^{2}\nonumber\\
&=&\frac{1}{C^{2}_{\varphi}}\Bigg|\iint_{\calG} (W_{\varphi}f)(s,t) \nonumber\\
&&\qquad\times \Bigg( (W_{\psi_1}\varphi)\left(\frac{a}{s},
\frac{b-t}{s}\right)
     -(W_{\psi_1}\varphi)\left(\frac{a_{j}}{s}, \frac{b_{j,k}-t}{s}\right)\Bigg)
  \frac{dsdt}{s^{d+1}}\Bigg|^2 \nonumber\\
&\le& \frac{1}{C^2_{\varphi}} \iint_{\calG} |(W_{\varphi}f)(s,t)|^2 \nonumber\\
&&\qquad \times\left| (W_{\psi_1}\varphi)\left(\frac{a}{s},
\frac{b-t}{s}\right)
    -(W_{\psi_1}\varphi)\left(\frac{a_{j}}{s}, \frac{b_{j,k}-t}{s}\right)\right|
   \frac{dsdt}{s^{d+1}}  \nonumber\\
&&\qquad \times   \iint_{\calG}\!
 \left|(W_{\psi_1}\varphi)\left(\frac{a}{s}, \frac{b-t}{s}\right)
    -(W_{\psi_1}\varphi)\left(\frac{a_{j}}{s}, \frac{b_{j,k}-t}{s}\right)\right|
 \frac{ ds dt}{s^{d+1}} \nonumber\\
&=& \frac{1}{C^2_{\varphi}} ....  \label{eq:a0}
\end{eqnarray}


\subsection{标点}
科技文献中一般用半角标点, 请参考《中国科学》发表的论文.

如果使用全角标点, 可以使用
\begin{verbatim}
  \punctstyle{<style>}
\end{verbatim}
选择标点样式, 有效值为
\begin{verbatim}
  quanjiao (所有标点符号占一个汉字宽度，
            相邻两个标点占一个半汉字宽度)
  banjiao  (所有标点符号占半个汉字宽度)
  hangmobanjiao (所有标点符号占一个汉字宽度，行末行首半角)
  kaiming  (句号、叹号、问号占一个汉字宽度，其他标点占半个汉字宽度)
\end{verbatim}
缺省为全角式。注意：不论选择哪种样式，都提供行末对齐(margin kerning)功能。



\begin{Theorem} \label{thm:latex}
\LaTeX 的输出是最完美的.
\end{Theorem}

先证明一个引理
\begin{Lemma} \label{thm:tex}
\TeX 文件在不同操作系统下的排版结果完全一致.
\end{Lemma}

\begin{proof}
这是证明.
\end{proof}


\begin{proof}[定理~\ref{thm:latex}的证明]
显然是错的.
\end{proof}

单个带编号的表达式
\begin{equation}\label{eq:a1}
x=y+z
\end{equation}

单个不带编号的表达式
\[
y=x-z.
\]

不带编号的多行表达式
\begin{eqnarray*}
x&=&y+z \\
 &=&z-s\\
 &<& 3. \\
 && \mbox{一些注释}
\end{eqnarray*}

带编号的多行表达式
\begin{eqnarray}
 x&=& y-z, \label{eq:aa1}\\
 y&=& x+z, \nonumber \\
 z&=&y-x. \label{eq:aa2}
\end{eqnarray}



引用:   定理\ref{thm:latex}的推论是什么呢？
方程式编号:  由(\ref{eq:a1})(\ref{eq:aa2})式.

\subsection{列举环境:  enumerate}

环境 enumerate 已经被改写，增加了一个可选参数[字符串], 用以控制所进。例如，
\begin{verbatim}
  \begin{enumerate}
    \item This is an example.
    \item This is an example.
      \begin{enumerate}
      \item This is an example.
      \item This is an example.
    \end{enumerate}
  \end{enumerate}
  \begin{enumerate}[Mn]% 字符串"Mn"的宽度为增加的缩进。
                       % 缺省值为 [M]
    \item This is an example.
    \item This is an example.
      \begin{enumerate}[Mnn]% 字符串"Mnn"的宽度为增加的缩进。
      \item This is an example.
      \item This is an example.
    \end{enumerate}
  \end{enumerate}
\end{verbatim}
的输出为
  \begin{enumerate}
    \item This is an example.
    \item This is an example.
      \begin{enumerate}
      \item This is an example.
      \item This is an example.
    \end{enumerate}
  \end{enumerate}
  \begin{enumerate}[Mn]% 字符串"Mn"的宽度为增加的所进。
                       % 缺省值为 [M]
    \item This is an example.
    \item This is an example.
      \begin{enumerate}[Mnn]% 字符串"Mnn"的宽度为增加的所进。
      \item This is an example.
      \item This is an example.
    \end{enumerate}
  \end{enumerate}

\clearpage

\section{\TeX\ 简介}

以下内容是 milksea@bbs.ctex.org 撰写的关于\TeX\ 的简单介绍。
注意这不是一个入门教程，不讲 \TeX\ 系统的配置安装，也不讲具体的 \LaTeX\ 代码。
这里仅仅试图以一些只言片语来解释：
进入这个门槛之前新手应该知道的注意事项，以及遇到问题以后该去如何解决问题。

\subsection{什么是Latex，我是否应该选择它}

\TeX\ 是最早由高德纳（Donald Knuth）教授创建的一门标记式宏语言，
用来排版科技文章，尤其擅长处理复杂的数学公式。\TeX\ 同时也是处理这一语言的排版软件。
\LaTeX\ 是 Leslie Lamport 在 \TeX\ 基础上按内容/格式分离和模块化等思想建立的一集 \TeX\ 上的格式。

\TeX\ 本身的领域是专业排版（即方正书版、InDesign 的领域），
但现在 TeX/LaTeX 也被广泛用于生成电子文档甚至幻灯片等，\TeX\ 语言的数学部分
偶尔也在其他一些地方使用。但注意 \TeX\ 并不适用于文书处理（MS Office 的领域，以前和现在都不是）。

选择使用 \TeX/\LaTeX\ 的理由包括：
\begin{itemize}
\item 免费软件；
\item 专业的排版效果；
\item 是事实上的专业数学排版标准；
\item 广泛的西文期刊接收甚或只接收 LaTeX 格式的投稿；
\item[] ……
\end{itemize}
不选择使用 \TeX/\LaTeX\ 的理由包括：
\begin{itemize}
\item 需要相当精力学习；
\item 图文混合排版能力弱；
\item 仅流行于数学、物理、计算机等领域；
\item 中文期刊的支持较差；
\item[] ……
\end{itemize}

请尽量清醒看待网上经常见到的关于 \TeX\ 与其他软件的优劣比较和口水战。在选择使用或离开之前，请先考虑
\TeX\ 的应用领域，想想它是否适合你的需要。

\def\AAAA{}

\subsection{我该用什么编辑器？}

编辑器功能有简有繁，特色不一，因人而易。
基本功能有语法高亮、方便编译预览就很好了，扩充功能和定制有无限的可能。
初学者强烈推荐免费的VSCode + TexLive。
VSCode支持多种插件，如LaTex Workshop插件支持LaTeX的编译，
Git Project Manager支持版本同步等。


\subsection{我该去哪里寻找答案？}

0、绝对的新手，先读完一本入门读物，了解基本的知识。

1、无论如何，先读文档！绝大部分问题都是文档可以解决的。

2、再利用 Google 搜索，利用(bbs.ctex.org)版面搜索。

3、清楚、聪明地提出你的问题。


\subsection{我应该看什么 LaTeX 读物？}

这不是一个容易回答的问题，因为有许多选择，也同样有许多不合适的选择。
这里只是选出一个比较好的答案。更多更详细的介绍可以在版面和网上寻找（注意时效）。

近两年 \TeX\ 的中文处理发展很快，目前没有哪本书在中文处理方面给出一个最新进展的合适综述，
因而下面的介绍也不主要考虑中文处理。

\begin{enumerate}
\item 我可以阅读英文

\begin{enumerate}
\item 我要迅速入门：ltxprimer.pdf (LaTeX Tutorials: A Primer, India TUG)
\item 我要系统学习：A Guide to LaTeX, 4th Edition, Addison-Wesley
      有机械工业出版社的影印版（《LaTeX实用教程》）
\item 我要深入学习：要读许多书和文档，TeXbook 是必读的
\item 还有呢？去读你使用的每一个宏包的说明文档
\item 还有许多专题文档，如讲数学公式、图形、表格、字体等
\end{enumerate}

\item 我更愿意阅读中文
\begin{enumerate}
\item 我要迅速入门：lnotes.pdf (LaTeX Notes, 1.20, Alpha Huang)
\item 我要系统学习：《LaTeX2ε 科技排版指南》，邓建松（电子版）
 如果不好找，看《LaTeX 入门与提高》第二版，陈志杰等
\item 我要深入学习：TeXbook0.pdf (特可爱原本，TeXbook 的中译，xianxian)
\item 还有呢？英语，绝大多数 TeX 资料还是英文的
\end{enumerate}
\end{enumerate}


\subsection{插图格式}

常见的图片格式都支持，如PDF、PNG、JPG格式的图形。



\subsection{作图}

目前已经有很多优秀的\LaTeX\ 作图宏包，如 pgf/Tikz 和 pstricks，
两者都具有强大的作图能力。

\chapter{说明}



本模板参照南开大学学位论文写作规范编写，
仅仅提供了论文的基本格式，包括章节标题和正文字体、字号等等的设置。



您自愿使用这个模板。
提供本模板的目的是为了给您的论文写作带来方便，然而，
作者不保证这个模板完全符合学校的要求，也不对由此产生的任何后果负责。
如果您不同意这些条款，请不要使用这个模板。


参考文献的录入请参考\ref{manual:ref}。 
%\include{references}


\bibliographystyle{NKThesis}
%\bibliographystyle{gbt7714-author-year}
\bibliography{nkthesis.bib}

% -*- coding: utf-8 -*-

%\makeschapterhead{致谢}
\chapter*{致谢}
感谢使用本模板。
%\include{appendices}
% -*- coding: utf-8 -*-


\chapter*{个人简历}

\noindent 姓名：张十三\\
出生日期：1903年6月1日\\


\myPara{教育背景：} 

\noindent
2012年9月-2016年7月\quad 南开大学\quad 某某学院\quad 某某学\quad 学士 \\ 
2012年9月-2016年7月\quad 南开大学\quad 某某学院\quad 某某学\quad 硕士 \\

%\setmainfont{Arial}
\myPara{硕士期间发表的学术论文：}
\begin{enumerate}
	\item \textbf{Yu-Huan Wu}*, Yun Liu*, Xin Zhan, and Ming-Ming Cheng. P2T: Pyramid Pooling Transformer for Scene Understanding. IEEE Transactions on Pattern Analysis and Machine Intelligence (TPAMI), 2022.  (SCI一区，CCF-A类期刊，影响因子24.314)
	\item Yun Liu*, \textbf{Yu-Huan Wu}*, Peisong Wen, Yujun Shi, Yu Qiu, and Ming-Ming Cheng. Leveraging Instance-, Image-and Dataset-Level Information for Weakly Supervised Instance Segmentation. IEEE Transactions on Pattern Analysis and Machine Intelligence (TPAMI), 2022. (SCI一区，CCF-A类期刊，影响因子24.314)
	\item \textbf{Yu-Huan Wu}, Yun Liu, Jun Xu, Jia-Wang Bian, Yu-Chao Gu, and Ming-Ming Cheng. MobileSal: Extremely Efficient RGB-D Salient Object Detection. IEEE Transactions on Pattern Analysis and Machine Intelligence (TPAMI), 2021.  (SCI一区，CCF-A类期刊，影响因子24.314)
	\item \textbf{Yu-Huan Wu}, Shang-Hua Gao, Jie Mei, Jun Xu, Deng-Ping Fan, Rong-Guo Zhang, and Ming-Ming Cheng. JCS: An explainable COVID-19 diagnosis system by joint classification and segmentation. IEEE Transactions on Image Processing (TIP), 2021. (SCI一区，CCF-A类期刊，影响因子11.041)
	\item \textbf{Yu-Huan Wu}*, Yun Liu*, Le Zhang, Ming-Ming Cheng, and Bo Ren. EDN: Salient object detection via extremely-downsampled network. IEEE Transactions on Image Processing (TIP), 2022. (SCI一区，CCF-A类期刊，影响因子11.041)
	\item \textbf{Yu-Huan Wu}, Yun Liu, Le Zhang, Wang Gao, and Ming-Ming Cheng. Regularized Densely-connected Pyramid Network for Salient Instance Segmentation. IEEE Transactions on Image Processing (TIP), 2021. (SCI一区，CCF-A类期刊，影响因子11.041)
	\item Yun Liu*, \textbf{Yu-Huan Wu}*, Yunfeng Ban, Huifang Wang, and Ming-Ming Cheng. Rethinking computer-aided tuberculosis diagnosis. IEEE Conference on Computer Vision and Pattern Recognition (CVPR), 2020. (CCF-A类会议，口头报告)
	\item 刘云, 陆承泽, 李仕杰, 张乐, \textbf{吴宇寰}, 程明明. 基于高效的多尺度特征提取的轻量级语义分割. 计算机学报, 2022. (中国卓越期刊，影响因子5.402)

\end{enumerate}

\myPara{研究生期间主要获得的荣誉奖励：}

\begin{enumerate}
	\item 南开十杰，2022年
	\item 国家奖学金，2022年
\end{enumerate}


\end{document}
